% TOP_PROBLEM: {"id": 3146, "title": "Half-integral flow polynomial values", "category_id": 3, "category": {"id": 3, "name": "graph_theory", "display_name": "Graph Theory", "description": "Problems involving graphs, networks, and their properties.", "slug": "graph-theory", "order_index": 3, "created_at": "2026-07-31T15:26:25.671Z"}, "status": "open", "problem_number": "OPG-348", "set_id": 12, "statusQualification": "Flow-polynomial counts are specified as finite-group (modular) flows, resolving the source shorthand “number of nowhere-zero k-flows.”"}
% FIRST_POSTED: 2026-10-01
% ATTEMPT_STATUS: unresolved
% UnsolvedMath ID 3146; source catalogue code OPG-348
% !TeX program = lualatex
% Research attempt by GPT-6 Astra Ultra; no claim of a new complete solution.
\documentclass[11pt,a4paper]{article}
\usepackage[margin=25mm]{geometry}
\usepackage{fontspec}
\setmainfont{lmroman10-regular.otf}[BoldFont=lmroman10-bold.otf,
 ItalicFont=lmroman10-italic.otf,BoldItalicFont=lmroman10-bolditalic.otf]
\usepackage{amsmath,amssymb,mathtools}
\usepackage{unicode-math}
\setmathfont{latinmodern-math.otf}
\AtBeginDocument{\let\setminus\smallsetminus}
\usepackage{xurl,hyperref}
\hypersetup{colorlinks=true,urlcolor=blue}
\setcounter{secnumdepth}{0}
\setlength{\parindent}{0pt}
\setlength{\parskip}{5pt}
\setlength{\emergencystretch}{3em}
\begin{document}
\section*{Half-integral flow polynomial values}\label{3146}
{\small UnsolvedMath ID 3146 · OPG-348 · Graph Theory · OpenGarden}
\subsection{Definitions and mathematical statement}
Let $G=(V,E)$ be a finite connected loopless undirected multigraph;
parallel edges are regarded as distinct. Assume $G$ has no bridge,
where a bridge is an edge whose deletion disconnects the graph.
Choose an orientation of its edges. For a positive integer $q$, a
nowhere-zero $\mathbb Z/q\mathbb Z$-flow assigns a nonzero residue
$\phi(e)$ to every edge so that the sum entering each vertex equals
the sum leaving it, with equality in $\mathbb Z/q\mathbb Z$.

The flow polynomial $\Phi(G,x)$ is characterized by counting these
flows at every positive integer $q$. A definition valid directly
for all real $x$ is
\[
 \Phi(G,x)=\sum_{A\subseteq E}(-1)^{|E|-|A|}
 x^{|A|-|V|+c(A)},
\]
where $c(A)$ is the number of connected components of the spanning
graph $(V,A)$, including isolated vertices. Each exponent is
nonnegative. Reversing an edge and negating its value gives a
bijection of flow sets, so the count is independent of orientation.

The conjecture asks whether every such $G$ satisfies
\[
 \Phi(G,11/2)>0.
\]
The argument $11/2$ is a polynomial evaluation, not the order of
a finite group. In particular, positivity at neighbouring integer
arguments does not by itself give the requested inequality.
The counting convention uses modular flows; counting bounded
integer-valued flows instead would define a different counting
function and would not specify this polynomial.
\subsection{Short English statement}
Is the flow polynomial of every bridgeless connected graph strictly positive at five and a half?
\subsection{Research attempt}
\paragraph{Scope and process.}
This is a GPT-6 Astra Ultra research attempt dated 1 October 2026, using
primary-source browsing and elementary mathematical checks. The canonical
definition above is retained. Results below are restricted results or
reductions, not a claim of a new solution to the entire catalogue question.
No formal proof assistant or external mathematical referee checked this text.


\paragraph{Literature and scope.}
Mohar's original problem page [R1] separates positivity at $11/2$ from
integer flow existence. Jacobsen and Salas [R2] construct real flow roots
above $5$, so a blanket assertion of positivity throughout $(5,\infty)$
cannot serve as a proof of this fixed-value question. Their result does
not by itself give a counterexample at $11/2$. We establish exact
positivity for graphs built from parallel-edge bundles by subdivisions
and one-vertex gluings, including bridgeless cacti and theta graphs.
All identities below concern the modular flow polynomial in the definition.

\paragraph{Two operations preserving a positive evaluation.}
Subdividing an edge does not change the flow polynomial. Orient the new
path consistently. Conservation at each new degree-two vertex forces
all its edge values to agree, and replacing the path by its single
value is a bijection of nowhere-zero flows. This holds for every
positive integer group order, hence as a polynomial identity.

If $G$ is formed by identifying one vertex of $G_1$ with one vertex
of $G_2$, then
\[
 \Phi(G,x)=\Phi(G_1,x)\Phi(G_2,x).
\]
For the converse direction that is sometimes missed, conservation at
all vertices of $G_1$ except the identification vertex already implies
conservation at that vertex: sum all the equations in $G_1$, so each
internal edge cancels. A flow on the union thus restricts to separate
flows on its summands, and arbitrary flows on the summands combine.
Again, counting at every positive integer gives the polynomial identity.
These arguments also handle repeated one-vertex gluings by induction.

\paragraph{Explicit parallel-edge computation.}
Let $B_s$ have two vertices and $s\ge2$ parallel edges, all oriented
from the first vertex to the second. For integer $q$, write $a_s(q)$
for the number of $s$-tuples of nonzero residues whose sum is zero.
Given the first $s-1$ entries, the last is forced and is nonzero unless
the preceding sum is zero. Therefore
\[
 a_1(q)=0,\qquad a_s(q)=(q-1)^{s-1}-a_{s-1}(q),
\]
and induction gives
\[
 \Phi(B_s,x)=\frac{(x-1)^s+(-1)^s(x-1)}{x}.
\]
The numerator vanishes at $x=0$, so this is a polynomial, despite its
rational-looking form. At $x=11/2$, put $t=9/2$. For even $s$ its
numerator is $t^s+t>0$. For odd $s\ge3$ it is
$t(t^{s-1}-1)>0$. This proves the desired strict inequality for every
bundle with at least two edges. The excluded case $s=1$ gives zero,
exactly as expected for a bridge.

Subdividing the edges of $B_2$ gives any cycle and recovers
$\Phi(C_m,x)=x-1$. A graph consisting of three internally disjoint
paths with common distinct endpoints is a subdivision of $B_3$, so
its polynomial is $(x-1)(x-2)$. More generally, subdivide any of the
bundles and glue the resulting graphs together successively at single
vertices. The product formula proves positivity throughout this family.
In particular, a connected bridgeless cactus with $b$ cycle blocks has
polynomial $(x-1)^b$. The isolated-vertex graph, if included in the
bridgeless convention, has polynomial $1$ and causes no exception.

\paragraph{Verification and attempted induction.}
The script \texttt{scripts/check\_attempt\_batch02\_algebra\_2026\_10\_01.py}
evaluates the definition's subset sum using exact rational arithmetic
at $11/2$. It checks $B_s$ for $2\le s\le8$, cycles of lengths three
through eight, and a theta graph with path lengths two, three and four,
against the formulas above. Parallel edges have separate subset indices;
no simplification to a simple graph is made.

For a general edge that is neither a loop nor a bridge, the recurrence
$\Phi(G,x)=\Phi(G/e,x)-\Phi(G\setminus e,x)$ suggests induction.
Knowing that the two terms on the right are positive does not determine
the sign of their difference. Neither subdivision nor one-vertex gluing
controls a general graph with a highly connected core. Likewise,
nonnegative counts at nearby integer values impose no sign condition
between those integers. These are precise failures of the two most
immediate extensions of the restricted proof.

\paragraph{Final status and first remaining gap.}
\textbf{UNRESOLVED.} Positivity is proved for the stated constructive
family. The missing ingredient is a comparison between contraction and
deletion values, or another sign-preserving reduction covering arbitrary
bridgeless cores. No such comparison is obtained here.

\subsection{Sources}
\begin{itemize}
\item[S1] ULAM.AI and the UnsolvedMath Contributors, supplied problems.json, record 3146 (OPG-348), OpenGarden collection. Catalogue identity and source transcription; CC BY 4.0.
\url{https://huggingface.co/datasets/ulamai/UnsolvedMath}.
\item[S2] Open Problem Garden, Half-integral flow polynomial values. Original problem and discussion; the mathematical formulation below is expanded from the supplied source record.
\url{http://www.openproblemgarden.org/op/half_integral_flow_polynomial_values}.

\item[R1] Bojan Mohar, \emph{Flow polynomial}, original research-problem
page, checked 1 October 2026.
\url{https://www.sfu.ca/~mohar/Problems/P0301FlowPolynomial.html}.
\item[R2] Jesper L. Jacobsen and Jes\'us Salas,
\emph{Is the five-flow conjecture almost false?}, Journal of
Combinatorial Theory, Series B 103 (2013), 532--565. Primary source
for real flow roots above five; abstract checked 1 October 2026.
\url{https://arxiv.org/abs/1009.4062}.

\end{itemize}
\end{document}
